\chapter{Qué habría que sancionar}
\label{cap:propuestas}

Este apéndice es el único del libro que propone texto normativo, y conviene decir de entrada
en qué carácter lo hace.

\begin{cautela}
\textbf{Primero, una corrección al modo de plantear el problema.} Es tentador cerrar diciendo
que lo que falta es voluntad política. Puede ser cierto y este libro no tiene con qué
afirmarlo: no entrevistó a ningún concejal, a ningún diputado ni a ningún senador, y atribuir
intención sin haber preguntado es exactamente lo que el capítulo \ref{cap:metodo} se prohíbe.

Lo que el relevamiento sí estableció es más modesto y más útil: \textbf{en los ocho puntos
que el capítulo \ref{cap:plan} enumera, nadie está obligado a nada}. No hay norma que mande
publicar el atlas, ni que exija fundar por qué una coordenada se publica y otra no, ni que
obligue a cruzar dos registros del mismo Estado. Donde no hay deber, no hace falta voluntad
para incumplir: alcanza con seguir haciendo lo que se viene haciendo.

\textbf{Por eso lo que sigue son deberes, no exhortaciones.} Cada propuesta convierte un
margen en una obligación con plazo y con consecuencia. Sólo una crea un órgano nuevo ---el Comité de Cuenca de la Ley 7---, pocas requieren presupuesto y se indica cuáles, y ninguna depende de que alguien cambie de
opinión sobre el fondo del asunto.

\textbf{Segundo, sobre el texto.} Los artículos que se transcriben son \textbf{borradores de
contenido, no de técnica legislativa}. Indican qué debe decir la norma, no cómo debe
redactarse. Quien los tome deberá pasarlos por asesoría legislativa: la numeración, las
remisiones, las derogaciones y las cláusulas transitorias son oficio propio y este libro no lo
tiene.
\end{cautela}

\section{Nivel municipal --- Concejo Deliberante}

Son cuatro ordenanzas, y más adelante se agregan una quinta y una sexta. Caen dentro de competencias que los municipios de La Caldera
y de Vaqueros ya ejercen, ninguna requiere autorización provincial y ninguna tiene costo más
allá de horas de trabajo administrativo.

\subsection{Ordenanza 1 --- Publicidad de los anexos y las capas}

\textbf{Responde a:} los ocho planos del atlas del Código se publican sólo como imágenes a escala municipal, sin las capas georreferenciadas de las que el propio Código dice provenir, y el capítulo \ref{cap:cot} tuvo que buscar dos de ellos, además, en un expediente judicial.

\begin{ficha}{Contenido propuesto}
\textbf{Art.\ 1.} Los anexos, planos, cuadros y gráficos que integran el Código de
Ordenamiento Territorial, conforme su artículo 12, deberán publicarse en el sitio web oficial
del Municipio, en formato de documento y \textbf{en sus capas georreferenciadas de origen}, en
un plazo de sesenta días.

\textbf{Art.\ 2.} La publicación comprende toda modificación posterior, dentro de los diez
días de sancionada.

\textbf{Art.\ 3.} \textbf{Ningún acto administrativo podrá fundarse en un anexo, plano o capa
que no se encuentre publicado} conforme al artículo 1. La invocación de un instrumento no
publicado constituye vicio en la motivación del acto.

\textbf{Art.\ 4.} Las capas se publicarán en formato abierto e interoperable, con su sistema
de referencia declarado.
\end{ficha}

\textbf{Por qué el artículo 3 es el que importa.} Sin él, la ordenanza es una buena intención
con plazo. Con él, publicar deja de ser una tarea pendiente y pasa a ser condición de validez
de lo que se apruebe.

\subsection{Ordenanza 2 --- El ancho de la franja, en metros}

\textbf{Responde a:} el artículo 43 declara no urbanizables ambas márgenes del río y no dice
cuántos metros. El capítulo \ref{cap:cot} lo buscó en el articulado, en la planilla del Anexo
VI y en el plano del Anexo II, y no está en ninguno.

\begin{ficha}{Contenido propuesto}
\textbf{Art.\ 1.} Establécese que la franja no urbanizable correspondiente a las zonas de
Parque Urbano (PU) del artículo 43 del Código de Ordenamiento Territorial tiene un ancho de
\underline{\hspace{1.5cm}} metros, medidos \textbf{desde la línea de ribera determinada} por
la autoridad hídrica provincial hacia el interior de cada margen.

\textbf{Art.\ 2.} En los tramos donde no exista línea de ribera determinada, la franja se
medirá desde el borde superior de la barranca, y la factibilidad de localización
\textbf{quedará supeditada} a la determinación previa de la línea de ribera.

\textbf{Art.\ 3.} El ancho establecido se incorporará a la planilla de indicadores del Anexo
VI y se graficará, acotado, en el plano del Anexo II.

\textbf{Art.\ 4.} La franja se inscribirá como restricción al dominio en las matrículas
alcanzadas, a cuyo efecto se comunicará a la Dirección General de Inmuebles.
\end{ficha}

\begin{cautela}
\textbf{El número va en blanco a propósito.} Este libro no tiene con qué proponer una cifra:
no hizo estudio hidrológico ni tiene las capas. Lo que sostiene es que \textbf{debe haber
una}, que debe estar en metros y que debe decir desde dónde se mide. Fijarla es tarea técnica,
no editorial.

El artículo 4 es el de mayor efecto práctico y el de mayor resistencia esperable: hoy la
franja no consta en el folio real, de modo que quien compra un lote no se entera.
\end{cautela}

\subsection{Ordenanza 3 --- Registro público de factibilidades}

\textbf{Responde a:} no existe nómina pública de factibilidades de localización otorgadas y
denegadas, y el capítulo \ref{cap:loteo} debió reconstruir el circuito loteo por loteo.

\begin{ficha}{Contenido propuesto}
\textbf{Art.\ 1.} Créase el \textbf{Registro Público de Factibilidades de Localización}, de
consulta libre y gratuita, publicado en el sitio web del Municipio y actualizado
mensualmente.

\textbf{Art.\ 2.} Cada asiento consignará: expediente, solicitante, matrícula, superficie,
zona del Anexo II, fecha de solicitud, fecha de resolución, \textbf{resultado y fundamento
sintético}, y organismos provinciales que tomaron intervención.

\textbf{Art.\ 3.} Se incorporarán los expedientes resueltos desde la sanción del Código de
Ordenamiento Territorial, en un plazo de ciento ochenta días.

\textbf{Art.\ 4.} La falta de asiento no afecta la validez de lo resuelto, pero
\textbf{habilita a cualquier vecino a requerir su incorporación} y hace responsable al
funcionario a cargo.
\end{ficha}

\subsection{Ordenanza 4 --- Publicidad del dictamen del artículo 12}

\textbf{Responde a:} el artículo 12 del Código dispone que, ante contradicción entre el texto
y los planos, prevalece lo que determine el Ejecutivo Municipal previo dictamen de la
Autoridad de Aplicación. No dice que ese dictamen deba publicarse \textit{(cap.\ \ref{cap:cot})}.

\begin{ficha}{Contenido propuesto}
\textbf{Art.\ 1.} El dictamen y el acto que resuelvan una contradicción en los términos del
artículo 12 del Código de Ordenamiento Territorial \textbf{se publicarán} en el sitio web y en
el registro del artículo 1 de la Ordenanza 3, dentro de los diez días.

\textbf{Art.\ 2.} El acto expresará el criterio aplicado y \textbf{será invocable como
antecedente} en casos análogos, debiendo el Ejecutivo fundar todo apartamiento.
\end{ficha}

\textbf{Es la más barata de las cuatro y la que más toca el fondo del asunto.} No manda
resolver de un modo ni de otro: manda decir con qué criterio se resolvió, y atarse a él la
próxima vez.

\section{Nivel provincial --- Cámara de Diputados y Senado}

Son seis leyes, a las que las secciones siguientes agregan cuatro más sobre el régimen del
agua, una sobre los regímenes promocionales, tres sobre el texto de los edictos y una sobre las
audiencias públicas viales. Cuatro no tienen costo. Dos lo tienen y se indica.

\subsection{Ley 1 --- Publicidad de las coordenadas de línea de ribera}

\textbf{Responde a:} dos determinaciones del mismo organismo sobre la misma matrícula,
separadas por cinco meses --- capítulo \ref{cap:ribera} ---: una publica ciento trece pares de
coordenadas y la otra las remite a la sede.

\begin{ficha}{Contenido propuesto}
\textbf{Art.\ 1.} Los edictos que publiquen determinaciones de línea de ribera deberán
contener \textbf{la totalidad de las coordenadas} de los puntos determinados, con indicación
del sistema de referencia empleado.

\textbf{Art.\ 2.} Queda prohibido remitir las coordenadas a consulta en la sede del organismo
como sustituto de su publicación.

\textbf{Art.\ 3.} La autoridad de aplicación publicará y mantendrá actualizado un
\textbf{geoportal de líneas de ribera determinadas}, de acceso libre, con las capas
correspondientes en formato abierto.

\textbf{Art.\ 4.} El incumplimiento del artículo 1 \textbf{suspende el plazo para recurrir}
previsto en el artículo 10 del Anexo I del Decreto 1.989/02, que comenzará a correr desde la
publicación completa.
\end{ficha}

\textbf{El artículo 4 es el corazón.} Hoy el afectado tiene diez días hábiles para apelar un
acto cuyas coordenadas debe ir a buscar a Avenida Bolivia 4650. Si el plazo no corre hasta que
el dato esté publicado, la omisión deja de ser gratuita.

\subsection{Ley 2 --- Publicidad registral de la línea de ribera}

\textbf{Responde a:} la línea de ribera se determina, se publica por dos días y no queda
asentada en la matrícula \textit{(cap.\ \ref{cap:ribera})}. Quien compra un lote no tiene cómo enterarse.

\begin{cautela}
\textbf{Una precisión indispensable, del mismo orden que la de la Ley 3.} La comunicación a la
Dirección General de Inmuebles \textbf{ya está ordenada}: el artículo 10 del Anexo I del Decreto
1.989/02 dispone que el acto que fija la línea de ribera «\textit{deberá ser registrado en el
libro de Catastros previstos en el Código de Agua y \textbf{comunicado a la Dirección de
Inmuebles para su registración y archivo}}». Y el artículo 126 del Código manda anotar las cotas
en el Libro de Catastro.

\textbf{De modo que el deber existe desde 2002 y el problema es otro:} que esa registración no
se traduce en nada que el comprador pueda ver. No consta en el certificado de dominio, no se
grafica en los planos de mensura y no impide visar un plano que la ignore.

Esta ley, entonces, \textbf{no crea la comunicación: la vuelve oponible}. Y agrega lo único que
falta de verdad, que es el artículo 4.
\end{cautela}

\begin{ficha}{Contenido propuesto}
\textbf{Art.\ 1.} La comunicación a la Dirección General de Inmuebles prevista en el artículo
10 del Anexo I del Decreto 1.989/02 \textbf{dará lugar a la anotación de la línea de ribera como
restricción al dominio} en las matrículas alcanzadas, dentro de los treinta días de recibida.

\textbf{Art.\ 2.} Los certificados e informes de dominio consignarán la existencia de la
anotación y su fecha.

\textbf{Art.\ 3.} Los planos de mensura, unión, división o loteo que comprendan inmuebles
ribereños deberán graficar la línea de ribera determinada y citar la resolución que la fija.
\textbf{Sin ese requisito no podrán ser visados.}

\textbf{Art.\ 4.} En los inmuebles ribereños sin línea de ribera determinada, la Dirección
General de Inmuebles consignará esa circunstancia en el certificado.
\end{ficha}

\begin{cautela}
El artículo 4 parece menor y no lo es. \textbf{Hace visible la ausencia}: convierte «no hay
determinación» en un dato que consta, en lugar de un silencio que el comprador interpreta como
que no hace falta.
\end{cautela}

\subsection{Ley 3 --- Línea de ribera previa a la aprobación de loteos}

\textbf{Responde a:} el capítulo \ref{cap:loteo} documenta un loteo aprobado sobre arroyos
cuya línea de ribera se determinó sin publicar coordenadas, y el capítulo \ref{cap:resistencias}
recoge que la propia autoridad minera informó en 2022 que hay loteos y asentamientos dentro de
las líneas de ribera delimitadas por la autoridad hídrica.

\begin{cautela}
\textbf{Una precisión que cambia el sentido de esta ley.} El requisito \textbf{ya se exige}: la
propia Secretaría de Recursos Hídricos declara, en la presentación que el capítulo
\ref{cap:ribera} transcribe, que para loteos limitados por un río «se exige en primera medida
la determinación de la línea de ribera».

De modo que esta propuesta \textbf{no crea un deber: le da rango legal a uno que hoy está en un decreto reglamentario ---el anexo del Decreto 1682/19 exige certificado de no inundabilidad y, para la preventa, línea de ribera cuando corresponda--- y en la práctica de un organismo}, y sobre todo lo vuelve verificable desde afuera. Hoy nadie puede
saber, mirando un expediente de aprobación, si el requisito se cumplió.
\end{cautela}

\begin{ficha}{Contenido propuesto}
\textbf{Art.\ 1.} La aprobación de todo fraccionamiento, loteo, club de campo o urbanización
que comprenda inmuebles ribereños \textbf{requerirá línea de ribera determinada y vigente}
sobre la totalidad de los cursos que lo atraviesan o lindan.

\textbf{Art.\ 2.} El certificado de no inundabilidad \textbf{no sustituye} el requisito del
artículo anterior, y \textbf{consignará el número de la resolución} que fija la línea de ribera
invocada, o la constancia de que no existe determinación.

\textbf{Art.\ 3.} La autoridad hídrica se expedirá en un plazo máximo de
\underline{\hspace{1.2cm}} días desde la solicitud. Vencido, el interesado podrá requerir
pronto despacho.

\textbf{Art.\ 4.} Los fraccionamientos aprobados con anterioridad conservarán su validez, pero
la autoridad hídrica determinará la línea de oficio, con comunicación al municipio y a la
Dirección General de Inmuebles, dentro de \underline{\hspace{1.2cm}} años.
\end{ficha}

\textbf{Costo: el del artículo 4.} Determinar de oficio las líneas faltantes del corredor
requiere comisiones técnicas y trabajo de campo. Es, junto con la digitalización de la Ley 5, una de las dos propuestas de esta sección que suponen una partida, y conviene decirlo sin rodeos.

\subsection{Ley 4 --- Correspondencia entre registro minero y evaluación ambiental}

\textbf{Responde a:} veinte concesiones de áridos figuran vigentes; seis con el informe de
impacto ambiental expresamente no aprobado, y nueve a la firma del Secretario, una desde marzo
de 2020 --- capítulo \ref{cap:resistencias}.

\begin{ficha}{Contenido propuesto}
\textbf{Art.\ 1.} La autoridad minera y la autoridad ambiental publicarán semestralmente un
\textbf{padrón único} de concesiones de extracción de áridos sobre cauces, que consigne para
cada una: expediente, concesionario, municipio, curso, estado del informe de impacto ambiental
y \textbf{fecha de su presentación}.

\textbf{Art.\ 2.} Toda concesión cuyo informe permanezca \textbf{sin resolución por más de
ciento ochenta días} desde su presentación será consignada como tal en el padrón, con
indicación del organismo en el que se encuentra.

\textbf{Art.\ 3.} La condición de «vigente» en el registro minero \textbf{no implica
habilitación ambiental} y así se consignará expresamente en el padrón.

\textbf{Art.\ 4.} Las actas de inspección se publicarán, con los datos personales
resguardados.
\end{ficha}

\textbf{Lo que esta ley no hace.} No cierra ninguna cantera, no revoca ninguna concesión y no
prejuzga sobre ninguna. Solamente impide que los dos registros puedan decir cosas distintas
sin que se note.

\subsection{Ley 5 --- Separatas y anexos del Boletín Oficial}

\textbf{Responde a:} la edición del Boletín que publica el decreto que rige los veinticinco actos de línea de ribera de este departamento trae una línea que remite a una separata, y la separata
no está digitalizada --- capítulo \ref{cap:ribera}.

\begin{ficha}{Contenido propuesto}
\textbf{Art.\ 1.} Toda separata y todo anexo del Boletín Oficial integran la publicación y
\textbf{se incorporarán al repositorio digital} junto con la edición que los anuncia.

\textbf{Art.\ 2.} El repositorio digital comprenderá la totalidad de las ediciones publicadas,
con buscador de texto completo y descarga libre.

\textbf{Art.\ 3.} Las separatas y anexos anteriores a la vigencia de esta ley se incorporarán
en un plazo de \underline{\hspace{1.2cm}} meses, priorizando las que contengan textos
normativos vigentes.

\textbf{Art.\ 4.} \textbf{No podrá tenerse por publicado} un acto cuyo texto se remita a una
separata o anexo no incorporado al repositorio.
\end{ficha}

\textbf{Costo: digitalización.} Es acotado y por única vez, y no requiere personal permanente.

\subsection{Ley 6 --- Paridad de los dos circuitos}

\textbf{Responde a:} la asimetría que el capítulo \ref{cap:tierrafiscal} documenta. El Estado
adquiere por decreto administrativo un inmueble que ocupa hace cuarenta y siete años; un
ocupante particular en la misma situación debe litigar.

\begin{ficha}{Contenido propuesto}
\textbf{Art.\ 1.} Los decretos que declaren operada la prescripción adquisitiva a favor de la
Provincia \textbf{se publicarán íntegramente}, con el plano de mensura que los sustenta.

\textbf{Art.\ 2.} Créase un \textbf{procedimiento administrativo de regularización dominial}
para ocupantes de inmuebles del departamento que acrediten posesión pública, pacífica e
ininterrumpida por el plazo legal, con los mismos requisitos probatorios que la Provincia
aplica en su propio beneficio.

\textbf{Art.\ 3.} La autoridad de aplicación publicará la nómina de solicitudes, su estado y
su resultado.

\textbf{Art.\ 4.} Toda denegatoria será \textbf{fundada y notificada por escrito}, con
indicación del requisito incumplido y del modo de subsanarlo.
\end{ficha}

\begin{cautela}
Esta es la propuesta más ambiciosa del apéndice y la que más excede lo que el relevamiento
documentó. El capítulo \ref{cap:tierrafiscal} establece la asimetría en los expedientes que
pudo ver; \textbf{no establece que un procedimiento administrativo sea la solución correcta},
ni evaluó sus riesgos --- entre otros, que un trámite más simple facilite regularizaciones que
la vía judicial hoy filtra.

Se incluye porque el problema está documentado y porque una propuesta de este apéndice debe
poder discutirse. \textbf{No porque este libro sepa que funcionaría.}
\end{cautela}

\section{Nivel provincial --- el régimen del agua}

Cuatro leyes más, que responden al capítulo \ref{cap:aguabaja}. Se agrupan aparte porque no
corrigen un margen de procedimiento sino una ausencia de fondo: \textbf{no hay nadie a cargo
de la cuenca}.

\subsection{Ley 7 --- Autoridad de cuenca del río La Caldera}

\textbf{Responde a:} la cuenca es una, la comparten dos municipios, ninguno administra el
conjunto del cauce, y la única figura metropolitana que la provincia construyó para este
corredor es la del transporte de pasajeros.

\begin{cautela}
\textbf{El Código ya prevé dos figuras parecidas, y ninguna alcanza.} Los artículos 198 a 200
permiten que los consorcios de usuarios de una cuenca se agrupen en \textbf{asociaciones de
segundo grado}; el artículo 6 crea un \textbf{Consejo Asesor Provincial del Agua} integrado por los usuarios del agua pública. El artículo fue vetado parcialmente por el Decreto 4.913/98, pero \textbf{el texto vigente que publica el Digesto conserva esa cláusula}, y el Consejo funciona: en diciembre de 2010 sesionaba el primer lunes de cada mes en la Secretaría de Recursos Hídricos. Su diseño, sin embargo, muestra el límite que esta propuesta quiere corregir. Lo integran los consorcios de usuarios de riego, agrupados por el Decreto 3774/09 en \textbf{seis zonas que no siguen las cuencas}: La Caldera comparte la del «Valle de Lerma» con la Capital, Rosario de Lerma, Chicoana, Cerrillos, La Viña y Guachipas, mientras General Güemes, aguas abajo del mismo Mojotoro, integra la del «Valle de Siancas». Sus dictámenes \textbf{no son vinculantes}. Y lo preside la misma autoridad a la que asesora: el secretario de Recursos Hídricos, con doble voto en caso de empate, y el coordinador general de Gestión Hídrica como vicepresidente, frente a siete representantes de los usuarios. Desde el reordenamiento de diciembre de 2025, además, el organismo que debía presidirlo ya no tiene rango de Secretaría (capítulo \ref{cap:ribera}; apéndice \ref{cap:normativa}).

\textbf{Y no son figuras teóricas: hay dos consorcios funcionando en este departamento}, el
del Loteo Valle Alegre y el de Valle Hermoso --- Coronado Cueto, con asambleas publicadas entre
2022 y 2026 \textit{(cap.\ \ref{cap:aguabaja})}. De modo que la herramienta existe, se usa y el
organismo la convoca.

Ninguna de las dos sirve, aun así, para lo que esta cuenca necesita. El consorcio
\textbf{agrupa usuarios de una toma común} y no incluye a los municipios, ni a la autoridad
minera, ni a las comunidades; y su competencia es la administración del recurso, no el dictamen
sobre obras y fraccionamientos. El Consejo Asesor, que existe, es provincial, se organiza por zonas de riego y asesora sobre política general.

\textbf{Pero su existencia sugiere un camino más corto que esta ley}, y sería deshonesto no
decirlo: los artículos 198 a 200 permiten que los consorcios de una cuenca se agrupen en
segundo grado, y para eso \textbf{no hace falta ninguna norma nueva}. Si los dos consorcios
existentes y los que puedan constituirse se agruparan, habría un interlocutor colectivo de los
usuarios de esta cuenca mañana mismo.

\textbf{Lo que falta es un órgano de cuenca con las autoridades adentro}, y eso el Código no lo
tiene. Esta ley lo crea aprovechando, donde puede, lo que ya existe: el artículo 2 incorpora a
los usuarios, que es exactamente a quienes el consorcio de segundo grado reúne.
\end{cautela}

\begin{ficha}{Contenido propuesto}
\textbf{Art.\ 1.} Créase el \textbf{Comité de Cuenca del río La Caldera}, con jurisdicción
sobre la cuenca delimitada en la cartografía oficial de la autoridad hídrica, desde las
nacientes hasta la confluencia con el río Mojotoro.

\textbf{Art.\ 2.} Lo integran: la autoridad hídrica provincial, que lo preside; las
autoridades ambiental y minera; los municipios de La Caldera y de Vaqueros; la prestadora del
servicio de agua; \textbf{un representante de las comunidades indígenas} con relevamiento
territorial en la cuenca; y \textbf{representantes de usuarios} --- riego, consumo humano y
extracción de áridos ---.

\textbf{Art.\ 3.} Competencias: aprobar el balance hídrico de la cuenca; dictaminar con
carácter previo y obligatorio sobre toda concesión, obra de captación, extracción sobre cauce
y fraccionamiento ribereño; y proponer el orden de prioridades de obra.

\textbf{Art.\ 4.} El dictamen previo \textbf{no es vinculante}, pero el acto que se aparte de
él deberá fundar el apartamiento y publicarse.

\textbf{Art.\ 5.} El Comité sesiona en el territorio de la cuenca, con \textbf{acta pública}
de cada reunión.

\textbf{Art.\ 6.} Ningún municipio ajeno a la cuenca ejercerá potestades de vigilancia o prohibición sobre sus cursos, salvo por acuerdo celebrado en el ámbito del Comité.
\end{ficha}

\begin{cautela}
\textbf{Sobre el artículo 4.} Un dictamen vinculante sería más fuerte y también más fácil de
bloquear: alcanzaría con no reunir el Comité. Un dictamen obligatorio pero no vinculante, con
deber de fundar el apartamiento, produce el efecto que este libro busca sin crear un poder de
veto: \textbf{obliga a decir por qué}.

\textbf{Sobre el artículo 6.} La Ley 1269 ya no rige: el Digesto de la Cámara de Diputados la lista entre las «leyes de objeto cumplido / plazo vencido / no vigentes». Por eso el artículo no sustituye nada; declara un principio. La ley de 1929 es, con todo, el mejor antecedente de lo que el artículo evita: la Legislatura resolvió un problema de la cuenca dándole a un municipio ajeno potestades sobre un tramo del río. Esa ley caducó, pero la autoridad de cuenca que no creó sigue sin existir.
\end{cautela}

\subsection{Ley 8 --- Balance hídrico público}

\textbf{Responde a:} este libro no pudo establecer cuánta agua sale de la cuenca ni adónde
va. Sabe que la planta del embalse se inauguró en 2021 con una producción de cuatro mil metros cúbicos por hora, que la prestadora declaraba en 2022 mil doscientos del Acueducto Norte y novecientos del Wierna para unos veintidós mil usuarios de la capital, y que las dos fuentes no dicen lo mismo sobre qué sistema opera \textit{(cap.\ \ref{cap:expropiacion})}; no sabe cuánto de eso resiste un balance.

\begin{cautela}
\textbf{Y acá la precisión es la más fuerte de todo el apéndice.} Medir no sólo es posible:
\textbf{está mandado desde 1999}.

El artículo 258 del Código de Aguas ordenó a la Autoridad establecer el aforo de partida de cada
fuente \textbf{dentro de los cinco años} de su vigencia, y renovarlo cada diez. El artículo 161
ordenó un inventario de aguas en el mismo plazo, \textbf{actualizado anualmente}. Y el artículo
259 prohíbe otorgar concesiones con dotación permanente antes de tener el aforo definitivo.

\textbf{Esta ley no inventa una obligación de medir. Le agrega la única pieza que le falta a la
que ya existe: que el resultado se publique.} Un inventario que se lleva y no se publica cumple
el artículo 161 y no sirve para nada de lo que este libro necesitó.
\end{cautela}

\begin{ficha}{Contenido propuesto}
\textbf{Art.\ 1.} La autoridad hídrica publicará \textbf{anualmente} el inventario de aguas
del artículo 161 y el balance hídrico de cada cuenca de la provincia, con: \textbf{el aforo de
partida del artículo 258 y sus renovaciones}; caudal medio y mínimo registrado por estación de
aforo;
volumen captado por cada obra; \textbf{destino del volumen captado, por localidad}; y volumen
concesionado, discriminado por uso.

\textbf{Art.\ 2.} Publicará y mantendrá actualizado el \textbf{registro de concesiones de agua
pública}, con titular, curso, caudal, uso y carácter permanente o eventual.

\textbf{Art.\ 3.} Las series de las estaciones de aforo serán de \textbf{acceso libre y
descarga abierta}, con su metodología y sus períodos de interrupción declarados.

\textbf{Art.\ 4.} Donde no exista estación de aforo en operación sobre un curso concesionado,
el balance lo consignará expresamente.
\end{ficha}

\textbf{El artículo 4 es el de siempre en este apéndice:} hace visible la ausencia. Un balance
que declara dónde no mide vale más que uno que estima en silencio.

\subsection{Ley 9 --- Compensación de las obras de captación}

\textbf{Responde a:} durante la construcción del acueducto de Campo Alegre se eliminó el
bosque ribereño. Hay carteles de un plan de forestación compensatoria y la población pidió
inventarios forestales en la costanera. El capítulo \ref{cap:expropiacion} lo documenta con
testimonios.

\begin{ficha}{Contenido propuesto}
\textbf{Art.\ 1.} Toda obra de captación, conducción o embalse sobre cursos de la provincia
deberá incluir, en su evaluación de impacto, \textbf{inventario forestal previo} del tramo
afectado y \textbf{plan de restauración del bosque ribereño}, con superficie, especies,
cronograma y responsable.

\textbf{Art.\ 2.} El plan se publicará y su cumplimiento se verificará por la autoridad
ambiental con \textbf{informes de avance públicos}, hasta su recepción definitiva.

\textbf{Art.\ 3.} Para las obras ya ejecutadas sin plan verificado, la autoridad ambiental
practicará \textbf{inventario del estado actual} del tramo afectado y determinará si
corresponde restauración, en un plazo de \underline{\hspace{1.2cm}} meses.

\textbf{Art.\ 4.} Los inventarios y sus resultados integrarán el balance hídrico del artículo
1 de la Ley 8.
\end{ficha}

\textbf{Por qué importa más de lo que parece.} El bosque ribereño no es paisaje: es la defensa
que el capítulo \ref{cap:defensas} muestra que este departamento viene reemplazando por defensas de rama, piedra y, desde 1980, hormigón. Restaurarlo es, en los términos que la propia sentencia del capítulo
\ref{cap:amparo} ordena, \textbf{infraestructura verde por sobre la gris}.

\subsection{Ley 10 --- Los cursos que no se registran}

\textbf{Responde a:} un arroyo que nace y muere dentro de una misma finca no se inscribe en el
registro de aguas públicas, porque su uso no es una concesión. En una cabecera de cuenca, esa
categoría no es marginal: es agua invisible.

\begin{ficha}{Contenido propuesto}
\textbf{Art.\ 1.} Los cursos de agua que nacen y mueren dentro de un mismo inmueble
\textbf{se inscribirán en un registro de cursos de dominio privado}, con su ubicación,
extensión y caudal estimado, a los solos efectos del balance hídrico de la cuenca.

\textbf{Art.\ 2.} La inscripción \textbf{no altera la titularidad ni el régimen de uso} de
esas aguas, ni genera obligación de concesión, canon ni autorización previa.

\textbf{Art.\ 3.} La declaración la formula el titular del inmueble al tramitar mensura,
fraccionamiento o factibilidad de localización.

\textbf{Art.\ 4.} La autoridad hídrica podrá relevarlos de oficio con fines de balance, sin
ingreso al inmueble y sobre cartografía e imágenes disponibles.
\end{ficha}

\begin{cautela}
\textbf{El artículo 2 es la propuesta entera.} Sin él, esto sería un avance sobre el dominio
privado y este libro no lo propone: no tiene con qué sostener que ese régimen deba cambiar, y
cambiarlo abriría una discusión constitucional que excede largamente lo que el relevamiento
documentó.

Lo que sí sostiene es que \textbf{no se puede planificar una cuenca cuya agua no se sabe
cuánta es}. Registrar la existencia no es regular el uso.
\end{cautela}

\subsection{La pregunta que este apéndice no resuelve}

\begin{cautela}
El capítulo \ref{cap:aguabaja} termina preguntando \textbf{qué le corresponde al territorio
que produce el agua}, y este apéndice no lo responde.

Sería fácil escribir una ley de canon de cuenca: un porcentaje de lo facturado por la
prestadora, girado al departamento de captación. Varias provincias tienen figuras parecidas y
sería la propuesta más vistosa del libro.

\textbf{No se incluye}, y conviene decir por qué. Este relevamiento no estableció cuánta agua
sale, ni cuánto cuesta producirla, ni qué recibe hoy el departamento por otras vías --- coparticipación,
obra pública, empleo de la prestadora ---. Sin esos tres números, una alícuota sería una cifra
inventada con apariencia de propuesta técnica.

\textbf{La Ley 8 es el paso previo, y por eso está.} Cuando el balance hídrico se publique, la
discusión sobre la compensación podrá darse con datos. Hoy solo podría darse con consignas, y
para eso este libro no hace falta.
\end{cautela}

\section{Nivel municipal --- la factibilidad de agua}

\subsection{Ordenanza 5 --- Factibilidad de agua previa al fraccionamiento}

\textbf{Responde a:} el capítulo \ref{cap:redes} documenta que los emprendimientos resolvieron
su abastecimiento por cuenta propia mientras las obras públicas quedaban proyectadas, y el
capítulo \ref{cap:aguabaja} muestra que esos dos movimientos conviven sin tocarse.

\begin{ficha}{Contenido propuesto}
\textbf{Art.\ 1.} La factibilidad de localización de todo fraccionamiento requerirá
\textbf{constancia de factibilidad de provisión de agua} emitida por la prestadora del
servicio, o, en su defecto, \textbf{copia de la concesión de agua pública} que ampare el
abastecimiento propio proyectado.

\textbf{Art.\ 2.} Ambas constancias se incorporarán al registro público de factibilidades de
la Ordenanza 3 y serán de consulta libre.

\textbf{Art.\ 3.} El abastecimiento por pozo propio deberá informar \textbf{caudal, profundidad
y destino}, y su registro se comunicará a la autoridad hídrica provincial.
\end{ficha}

\section{Nivel provincial --- los regímenes promocionales}

Una sola ley, y antes conviene explicar por qué una sola.

\begin{cautela}
\textbf{Por qué este apéndice propone poco sobre el capítulo \ref{cap:trabajo}.}

La tentación es evidente. Un departamento que perdió el ochenta y ocho por ciento de su
superficie agropecuaria en dieciséis años parece pedir a gritos un régimen de promoción
agropecuaria, una línea de crédito, una exención simétrica a la que recibió el turismo.

\textbf{Este libro no tiene con qué proponer eso.} No estableció por qué desapareció esa
superficie ---el capítulo \ref{cap:trabajo} dice expresamente que no lo sabe---, no sabe
cuánto empleo generan los proyectos promovidos, y no evaluó ningún instrumento de fomento. Proponer política
productiva sobre esa base sería inventar, y la única diferencia con las propuestas anteriores
sería que nadie lo notaría.

Lo que sí puede proponerse es lo de siempre, y alcanza para bastante: \textbf{que lo que el
Estado promueve se sepa, se mida y se verifique}. Si después de eso alguien quiere discutir
política productiva, va a poder hacerlo con números.
\end{cautela}

\subsection{Ley 11 --- Publicidad y verificación de los regímenes promocionales}

\textbf{Responde a:} el Decreto 3069/14 ratifica un contrato de promoción turística cuyo
contenido no se publicó, y el Decreto 463/15 otorga exención de tributos y créditos fiscales
sin que conste el monto del beneficio. El capítulo \ref{cap:opacidad} lo
registra como caso y el capítulo \ref{cap:trabajo} muestra lo que significa: no se sabe cuánto
dinero público se resignó ni a cambio de qué.

\begin{ficha}{Contenido propuesto}
\textbf{Art.\ 1.} Los contratos, convenios y actos que otorguen \textbf{exenciones,
diferimientos o beneficios impositivos} se publicarán \textbf{íntegramente} junto con el acto
que los aprueba, con indicación de beneficiario, tributos alcanzados, plazo, monto estimado
del beneficio y \textbf{contraprestaciones comprometidas}.

\textbf{Art.\ 2.} El beneficiario presentará \textbf{informe anual} de cumplimiento de las
contraprestaciones, que incluirá \textbf{inversión ejecutada y puestos de trabajo declarados},
y que será público.

\textbf{Art.\ 3.} La autoridad de aplicación publicará anualmente el \textbf{gasto tributario}
resultante, \textbf{desagregado por régimen, por departamento y por beneficiario}.

\textbf{Art.\ 4.} La falta de publicación conforme al artículo 1 \textbf{suspende el goce del
beneficio} hasta que se subsane.

\textbf{Art.\ 5.} Los contratos vigentes a la sanción de esta ley se publicarán dentro de los
noventa días.
\end{ficha}

\textbf{El artículo 3 es el que convierte esto en un instrumento y no en un trámite.} Un
gasto tributario desagregado por departamento permite, por primera vez, poner al lado dos
cifras que hoy no existen: lo que la provincia resignó para promover una actividad en un
territorio y lo que invirtió en ese mismo territorio por otras vías.

\textbf{Y el artículo 2 responde a una carencia propia de este libro.} El capítulo
\ref{cap:trabajo} sabe cuánta gente trabaja en cada rama, pero no cuánto empleo generan los proyectos promovidos, porque
nadie lo publica. Si el beneficiario de una exención debe declarar sus puestos de trabajo, el
dato aparece sin necesidad de censo.

\section{Nivel municipal --- el registro de alojamientos}

\subsection{Ordenanza 6 --- Registro de alojamientos turísticos}

\textbf{Responde a:} el capítulo \ref{cap:poblacion} establece que el departamento tiene una
proporción de viviendas sin residentes permanentes que duplica la provincial, y no puede
distinguir cuánto de eso es vivienda de fin de semana y cuánto alojamiento turístico en
actividad.

\begin{ficha}{Contenido propuesto}
\textbf{Art.\ 1.} Créase el \textbf{Registro de Alojamientos Turísticos} del municipio, de
consulta pública, con denominación, domicilio, cantidad de plazas, modalidad y fecha de
habilitación.

\textbf{Art.\ 2.} Se incluirán las viviendas ofrecidas en alquiler temporario a través de
plataformas, con los mismos datos.

\textbf{Art.\ 3.} El registro se publicará semestralmente en forma agregada por barrio o
paraje, \textbf{resguardando los datos personales de los titulares}.

\textbf{Art.\ 4.} La información se remitirá al organismo provincial competente en turismo y
al que lleva las estadísticas del municipio.
\end{ficha}

\begin{cautela}
El artículo 3 no es un detalle. Un registro de alojamientos publicado con nombre y domicilio
de cada titular es, en un departamento de doce mil habitantes, \textbf{un padrón de
quién tiene una casa vacía}. Agregado por barrio sirve para planificar; individualizado sirve
para otras cosas.
\end{cautela}

\section{Nivel provincial --- lo que el propio edicto debería decir}

Tres leyes más, y las tres salen de una misma constatación: \textbf{los defectos que este libro
encontró no están en las decisiones, están en los textos que las publican}.

\subsection{Ley 12 --- Correspondencia entre el aviso de inicio y la determinación}

\textbf{Responde a:} el hallazgo central del capítulo \ref{cap:ribera}. Dentro del expediente
34-167.127/13, la Resolución 395/13 aprueba la comisión para el río La Caldera y el arroyo El
Manzano; la determinación de julio de 2014 recae sobre esos dos cursos y publica sus
coordenadas; y la de diciembre de 2014 recae sobre \textbf{otros dos arroyos que ningún acto de
comisión nombra}.

\begin{ficha}{Contenido propuesto}
\textbf{Art.\ 1.} El aviso de inicio del artículo 4 del Anexo I del Decreto 1.989/02 deberá
\textbf{individualizar los cursos de agua} sobre los que la comisión técnica realizará la
determinación, además del inmueble.

\textbf{Art.\ 2.} \textbf{La determinación no podrá exceder los cursos anunciados.} La
incorporación de un curso no comprendido en el aviso requiere \textbf{nuevo aviso} y reabre el
plazo del artículo 5 respecto de ese curso.

\textbf{Art.\ 3.} El edicto de determinación \textbf{citará el acto de comisión} que le dio
origen, con su número y fecha de publicación.

\textbf{Art.\ 4.} La falta de la cita del artículo 3 no invalida el acto, pero \textbf{suspende
el plazo para recurrir} hasta que se subsane.
\end{ficha}

\textbf{Es la propuesta más directamente derivada de este libro}, y la única que habría hecho
visible, en su momento y sin necesidad de este relevamiento, lo que el capítulo \ref{cap:ribera} tardó decenas de páginas en establecer: que dos arroyos entraron a una
determinación sin haber sido anunciados.

\subsection{Ley 13 --- Clasificación de los avisos oficiales}

\textbf{Responde a:} determinaciones de línea de ribera del mismo departamento aparecen, según
el año, bajo \textbf{«determinaciones de línea de ribera»}, \textbf{«avisos administrativos»} o
\textbf{«concesiones de agua pública»}. El capítulo \ref{cap:ribera} encontró los actos de 2009
y 2011 sólo porque buscó por texto.

\begin{ficha}{Contenido propuesto}
\textbf{Art.\ 1.} El Boletín Oficial publicará un \textbf{nomenclador de tipos de aviso}, con
la definición de cada uno y los actos que comprende.

\textbf{Art.\ 2.} Cada aviso se clasificará conforme a ese nomenclador, y el organismo
requirente indicará el tipo al solicitar la publicación.

\textbf{Art.\ 3.} El buscador público permitirá filtrar por \textbf{tipo, organismo,
departamento y rango de fechas}, y \textbf{buscar por texto completo} en el cuerpo del aviso y
no sólo en su título.

\textbf{Art.\ 4.} Los avisos ya publicados se reclasificarán progresivamente, \textbf{sin
alterar su texto}.
\end{ficha}

\begin{cautela}
\textbf{Es de las propuestas más baratas del apéndice y quizá la de mayor alcance.} No cambia nada de
lo que los organismos deciden ni de lo que publican: cambia si eso puede encontrarse.

Este libro existe, en buena medida, porque alguien se sentó a buscar en un archivo mal
indexado. \textbf{Un índice mejor no habría cambiado un solo hecho, pero habría hecho
innecesarias varias de sus páginas}, y eso es exactamente lo que una
infraestructura de publicidad debería lograr.
\end{cautela}

\subsection{Ley 14 --- Fe de erratas y control del texto publicado}

\textbf{Responde a:} tres errores de texto que el capítulo \ref{cap:ribera} documenta, todos arrastrados de un
edicto al siguiente. Un decreto citado como «1898/02» en edictos de distintos años entre 2013 y 2025 ---no en todos---, cuando es el
1.989/02. Un edicto de 2026 que, firmado por la \textbf{Sub}secretaría, ordena recurrir «ante la
Secretaría». Y un edicto de 2011 que, sobre un inmueble lindero al arroyo La Calderilla, dice
que la comisión propone la línea de ribera \textbf{«de ambas márgenes del río Arenales»}, que
es el río de la ciudad de Salta.

\begin{ficha}{Contenido propuesto}
\textbf{Art.\ 1.} El organismo requirente \textbf{verificará y certificará}, al solicitar la
publicación, que el texto del aviso concuerda con el acto en cuanto a \textbf{número de norma,
identificación del inmueble y denominación del curso o bien de que se trate}.

\textbf{Art.\ 2.} Advertida una discordancia, el organismo publicará \textbf{fe de erratas}
dentro de los diez días, con el mismo alcance y por el mismo término que la publicación
original.

\textbf{Art.\ 3.} La fe de erratas \textbf{reabre el plazo para recurrir} cuando la
discordancia recaiga sobre alguno de los elementos del artículo 1.

\textbf{Art.\ 4.} Cualquier persona podrá denunciar la discordancia, y la denuncia se resolverá
por acto fundado y publicado.
\end{ficha}

\begin{cautela}
\textbf{Sobre el alcance de esto, y para que no se lea como no se debe.} Ninguno de los tres
errores que este libro encontró parece haber causado daño a nadie, y ninguno de los tres
justifica, por sí mismo, cuestionar el acto que lo contiene. Un error de tipeo en un edicto es
un error de tipeo.

\textbf{Lo que los vuelve dignos de una norma es su persistencia.} No son tres descuidos
distintos: son textos que pasaron de una publicación a otra sin que nadie los leyera, y en el
caso del decreto mal numerado, durante años. Y en el caso del río equivocado, el texto arrastrado está en el
instrumento que fija, sobre un inmueble concreto, \textbf{dónde termina la propiedad privada y
empieza el dominio público}.
\end{cautela}

\subsection{Ley 15 --- Accesibilidad de las audiencias públicas de proyectos viales}

\textbf{Responde a:} la audiencia pública del Estudio de Impacto Ambiental y Social de la obra
de pavimentación de la Ruta 9 entre la capital y La Caldera se celebró \textbf{en el municipio
de Vaqueros}, un lunes a las nueve de la mañana; el plazo para constituirse en parte fue de \textbf{ocho días} desde la primera publicación; y el expediente sólo podía consultarse
\textbf{en la ciudad de Salta, de nueve a trece} \textit{(cap.\ \ref{cap:vaqueros})}.

\begin{ficha}{Contenido propuesto}
\textbf{Art.\ 1.} Cuando una obra vial atraviese o afecte a más de un municipio, la audiencia
pública \textbf{se celebrará en cada uno de ellos}, o en su defecto en el de mayor distancia a
la sede del organismo, y \textbf{fuera del horario laboral} o en día no laborable.

\textbf{Art.\ 2.} El plazo para constituirse en parte no será menor a \textbf{veinte días
hábiles} contados desde la última publicación.

\textbf{Art.\ 3.} El Estudio de Impacto Ambiental y Social estará \textbf{disponible en línea}
desde la primera publicación de la convocatoria, y en soporte papel \textbf{en la sede del
municipio o los municipios afectados}.

\textbf{Art.\ 4.} El acta de la audiencia, las presentaciones recibidas y \textbf{la respuesta
fundada del organismo a cada observación} se publicarán dentro de los treinta días.

\textbf{Art.\ 5.} El incumplimiento de los artículos 1 a 3 \textbf{suspende el plazo} para
impugnar el acto que apruebe el estudio.
\end{ficha}

\begin{cautela}
\textbf{Una precisión sobre lo que esta ley no hace}, del mismo orden que las anteriores. El
artículo 49 de la Ley 7070 ya exige la audiencia pública, y el Reglamento de Audiencia Pública
de Proyectos Viales de la Dirección de Vialidad ya regula su trámite. \textbf{Nada de lo
propuesto crea la audiencia: sólo fija dónde, cuándo y con qué anticipación.}

\textbf{Y conviene decir lo que este libro no sabe.} No estableció que la audiencia de noviembre
de 2023 se haya celebrado mal, ni que alguien haya quedado excluido: no tiene el acta. Lo que
registra es que \textbf{las condiciones de acceso hacían improbable la participación de quien
viviera en el municipio del extremo de la obra}, y eso puede corregirse sin discutir ningún
caso.
\end{cautela}

\section{Nivel ejecutivo --- decretos y resoluciones}

Cinco actos que no requieren ley y que podrían dictarse esta semana.

\begin{enumerate}[leftmargin=*,itemsep=0.5em]
\item \textbf{Decreto que ordena publicar el Libro de Catastro de Aguas.} El Anexo I del
Decreto 1.989/02 manda registrar allí cada determinación. Publicarlo no requiere norma nueva:
requiere decidirlo.

\item \textbf{Resolución conjunta de las Secretarías de Minería y de Ambiente} que instruya el
cruce de padrones previsto en la Ley 4 de este apéndice. Lo que la ley haría obligatorio y
permanente, una resolución puede hacerlo ahora y por una vez.

\item \textbf{Resolución de la autoridad hídrica que fije y publique el criterio} de
publicación de coordenadas, mientras no exista la Ley 1. Basta con que diga qué se publica
siempre y por qué. Es el acto más barato de todo este apéndice y el que más directamente toca
lo que el capítulo \ref{cap:opacidad} describe.

\item \textbf{Resolución que publique las series de las estaciones de aforo} de la cuenca y
declare cuáles están fuera de servicio. El capítulo \ref{cap:aguabaja} muestra que este libro
no pudo establecer cuánta agua baja; publicar los caudales máximos instantáneos que el SNIH no ofrece y declarar formalmente que hoy ninguna estación opera dentro de la cuenca (capítulo \ref{cap:ribera}) no requiere ley, no requiere
presupuesto y \textbf{es la condición de cualquier planificación posterior}.

\item \textbf{Publicación del gasto tributario por departamento} correspondiente a los
regímenes promocionales vigentes. Es el artículo 3 de la Ley 11 de este apéndice, y no
requiere esperarla: la información existe en la administración tributaria y lo único que falta
es desagregarla y publicarla.
\end{enumerate}

\section{Cómo se presenta cada una}

\begin{itemize}[leftmargin=*,itemsep=0.35em]
\item Las \textbf{ordenanzas} las presenta cualquier concejal de La Caldera o de Vaqueros, o
el Departamento Ejecutivo Municipal. También puede hacerlo un vecino, por nota, para que el
Concejo la tome como proyecto.
\item Las \textbf{leyes} las presenta cualquier diputado o senador provincial. El departamento
La Caldera elige un senador y un diputado propios.
\item Los \textbf{decretos y resoluciones} no requieren presentación: se dictan.
\end{itemize}

\textbf{Y conviene decir algo sobre cómo se usa este libro en esa gestión.} Sirve para mostrar
que el problema está documentado, con expediente y fecha. \textbf{No sirve como fundamento de
una acusación}, y quien lo use de ese modo lo estará usando contra su propio método.

\begin{cautela}
Una observación final, que vale para las veintiséis propuestas.

\textbf{Ninguna manda decidir de un modo determinado.} No fijan el ancho de la franja, no
dicen qué canteras deben cerrarse, no resuelven si un loteo estuvo bien aprobado. Mandan
\textbf{publicar, fundar, registrar y cruzar}.

Esa elección es deliberada y es la conclusión de todo el libro. Un territorio no se ordena
prohibiendo más: se ordena haciendo que cada decisión deje rastro y tenga que explicarse.
\textbf{Lo que este relevamiento encontró no fue un ordenamiento ausente, sino uno que
funciona sin tener que dar razones.}
\end{cautela}
